\section{Energy and free-energy estimates}\label{sec:local-energy}\label{sec:energy}

This section gives the particle-number costs needed to localize the gas.
Throughout, $0<T\leq\rho^2$, $\beta=T^{-1}$, and
$\rho/2\leq N/V\leq2\rho$. All statements hold after decreasing a
potential-dependent upper bound on $\rho$ and increasing the volume
threshold at each fixed density. The Hamiltonian with a bounded external
potential $-W$ is denoted by $H_{j,L}^{(-W)}$; its symmetric partition
function and free energy are $Z_j^{(-W)}$ and
$F_j^{(-W)}=-T\log Z_j^{(-W)}$. We set $Z_0=1$ and $F_0=0$.
For a trace-one state $\sigma$, write
$\mathcal F^{(-W)}(\sigma)=\operatorname{Tr}H^{(-W)}\sigma-TS(\sigma)$.
Here $S(\sigma)=-\operatorname{Tr}(\sigma\log\sigma)$ is the
von Neumann entropy, with $0\log0=0$.
Constants may depend on $v$, and later on already fixed dimensionless
parameters, but never on $\rho,L,N,T$.

\begin{lemma}[A Neumann cell bound]\label{lem:neumann-cell}
There are $b_0,c_v>0$, depending only on the fixed potential, such that
for a cube $Q$ of side $b\geq b_0$, every $y\in Q$, and every
$f\in H^1(Q)$,
\begin{equation}
 \int_Q\bigl(|\nabla f|^2+v(x-y)|f|^2\bigr)\,dx
 \geq c_v b^{-3}\int_Q|f|^2\,dx.\tag{1}\label{eq:1}
\end{equation}
The same conclusion holds for fixed positive coefficients in front of
the two terms, with changed constants. Consequently, for disjoint
cubes $Q$ of common side $b$, kinetic energy plus nonnegative pair
energy bounds below the multiplication operator
\begin{equation}
 c_v b^{-3}\sum_Q n_Q\mathbf1_{\{n_Q\geq2\}}.\tag{2}\label{eq:2}
\end{equation}
Both statements apply to torus cubes smaller than the torus, and to
Dirichlet states extended by zero.
\end{lemma}

\begin{proof}
Choose $\nu>0$ for which the radial set
$A=\{x:v(x)\geq\nu\}$ has positive measure. Increase $b_0$ so that
$b_0/2>R_v$ and $b_0\geq1$. For each $y\in Q$ choose, independently
in the three coordinate directions, an orientation with available
length at least $b/2$. The resulting octant of $y+A$ lies in $Q$ and
has measure $|A|/8=:a_v>0$ by radial symmetry. Call it $A_y$.
For $f_0=|Q|^{-1}\int_Q f$, the scale-invariant Sobolev--Poincar\'e
inequality gives $\|f-f_0\|_6\leq C\|\nabla f\|_2$. Thus
\[
 a_v|f_0|^2\leq2\int_{A_y}|f|^2+
 2\int_{A_y}|f-f_0|^2
 \leq2\nu^{-1}\int_Qv(x-y)|f|^2+
 C a_v^{2/3}\int_Q|\nabla f|^2.
\]
Moreover $\|f-f_0\|_2^2\leq Cb^2\|\nabla f\|_2^2$.
Since $b^2\leq b^3$ for $b\geq1$, these estimates give (1).
Positive fixed energy coefficients only change the constants.

For (2), cut configuration space according to the cube containing
each coordinate and discard cross-cube interactions. This is Neumann
bracketing: the restrictions of an $H^1$ function to the pieces are
admissible with no boundary condition at the cuts. Within a cube
containing $j\geq2$ coordinates, pair them disjointly. For each pair
apply (1) in the first coordinate with the other coordinate held
fixed, and then integrate the other coordinates. No coordinate
kinetic term or pair interaction is used twice. The number of pairs
is $\lfloor j/2\rfloor\geq j/3$. Summing the resulting inequalities
proves (2). The argument is unaffected by additional nonnegative
periodic images of $v$. Zero extension gives the Dirichlet assertion.
\end{proof}

\begin{lemma}[An entropy bound]\label{lem:kinetic-entropy}
If $\Omega$ is an open subset of the torus with Dirichlet kinetic
energy, or the full torus with periodic kinetic energy, then for
every $\alpha>0$ and every finite-energy symmetric trace-one $j$-particle state,
\begin{equation}
 \alpha\operatorname{Tr}\Bigl(\sum_{i=1}^j-\Delta_i\Bigr)\sigma
 -TS(\sigma)
 \geq-Tj-C_\alpha T|\Omega|(V^{-1}+T^{3/2}).\tag{3}\label{eq:3}
\end{equation}
In particular a state of finite kinetic energy has finite entropy.
\end{lemma}

\begin{proof}
The periodic heat kernel satisfies
$K_t(0)\leq C(V^{-1}+t^{-3/2})$ by its Gaussian image sum,
or equivalently by its Fourier series. Killing at $\partial\Omega$
decreases its diagonal, so
$\operatorname{Tr}e^{t\alpha\Delta_\Omega}
\leq C_\alpha|\Omega|(V^{-1}+t^{-3/2})$.
The ideal Bose grand partition function with fugacity $e^{-1}$ has
logarithm
\[
 \sum_{l\geq1}\frac{e^{-l}}l
 \operatorname{Tr}e^{l\beta\alpha\Delta_\Omega}
 \leq C_\alpha|\Omega|(V^{-1}+T^{3/2}).
\]
Its $j$-particle term is $e^{-j}$ times the canonical partition
function. Applying the Gibbs variational inequality to that term
gives (3). For an arbitrary finite-energy state, the same inequality
follows first on finite spectral truncations and then by monotone
entropy approximation, or directly from nonnegativity of relative
entropy to the ideal Gibbs state. The right side is finite, proving
the entropy assertion as well.
\end{proof}

\begin{proposition}[Particle-number free-energy bounds]\label{prop:increments}
Uniformly in the temperature range under consideration,
\begin{align}
 F_N&\geq c_v\rho N,\tag{4}\label{eq:4}\\
 F_{N+k}&\geq F_N+c_v\rho k\qquad(k\geq0).\tag{5}\label{eq:5}
\end{align}
For every bounded $W\geq0$, with no upper bound on its size required,
\begin{equation}
 F_j^{(-W)}-F_{j-1}^{(-W)}\leq C_v\rho
 \qquad(1\leq j\leq4\rho V).\tag{6}\label{eq:6}
\end{equation}
The constant in (6) is independent of $W$.
\end{proposition}

\begin{proof}
Partition the torus into equal cubes of side in $[b,2b]$, where
$b=(4/\rho)^{1/3}$. The number of cubes is at most
$V/b^3=\rho V/4\leq N/2$ when $L$ is sufficiently large.
Reserve a fixed fraction of kinetic energy for (3), and apply (2)
to the rest of the kinetic energy and the interaction. Since
$\sum_Q n_Q\mathbf1_{\{n_Q\geq2\}}\geq N-\#\{Q\}\geq N/2$,
the energy contribution is at least $c_v\rho N$.
The entropy loss is bounded by
$TN+C T(1+VT^{3/2})$. Divided by $\rho N$, these terms tend uniformly
to zero as $\rho\downarrow0$ and the volume threshold increases:
$T/\rho\leq\rho$, $T/(\rho N)\leq\rho/N$, and
$VT^{5/2}/(\rho N)\leq2\rho^3$. This proves (4), after changing $c_v$.

We next show that $F_j/j$ is nondecreasing in $j\geq1$.
Regard a symmetric $j$-particle state as a state on the labeled tensor
product and denote by $S_i$ the entropy of its trace-one
$i$-particle marginal; put $S_0=0$. Strong subadditivity~\cite[Theorem~2 and Section~IV]{LiebRuskai1973Entropy} applied to
two overlapping $i$-coordinate subsystems with overlap of size
$i-1$ gives $2S_i\geq S_{i-1}+S_{i+1}$. The case $i=1$ uses
ordinary subadditivity. Permutation symmetry identifies the
marginals occurring here. Therefore $S_i$ is concave and
$S_j/j\leq S_{j-1}/(j-1)$ for $j\geq2$.
The one-body energy per particle is unchanged by passage to the
$(j-1)$-particle marginal. The pair energy per particle decreases,
because its coefficient is $(j-1)/2$ before passage and $(j-2)/2$
afterwards, and the common two-body expectation is nonnegative;
for $j=2$ the latter energy is zero. The marginal remains supported
on the symmetric subspace. Applying these observations to the Gibbs
state and then using its marginal as a variational state proves
$F_j/j\geq F_{j-1}/(j-1)$. All entropies are finite by (3).
Consequently
$F_{N+k}\geq(N+k)F_N/N$, which with (4) proves (5).

The lower increments now quantify the free-energy gain from adding
particles. To obtain the complementary upper cost, uniformly even in
an attractive external field, we insert one particle directly in the
path representation.

Retain in the symmetrized Feynman--Kac
formula only those permutations in which the extra label is fixed.
There are $(j-1)!$ such permutations against the factor $1/j!$.
The added particle is a free torus loop whose start is integrated
over the torus. After normalization by $VK_\beta(0)$ its position
at every time is uniform, independently of all old trajectories.
Thus its expected total interaction with the old paths is
$\beta(j-1)\|v\|_1/V$. Its external weight is at least one because
$W\geq0$. Jensen's inequality, conditional on all old paths, gives
\[
 Z_j^{(-W)}\geq\frac{VK_\beta(0)}j
 e^{-\beta(j-1)\|v\|_1/V}Z_{j-1}^{(-W)}.
\]
Since $K_\beta(0)\geq(4\pi\beta)^{-3/2}$, the positive part of the
prefactor's contribution to the increment is bounded by
$T\log_+(C\rho T^{-3/2})$. For small $\rho$, the derivative of
$t\mapsto t\log(C\rho t^{-3/2})$ on $0<t\leq\rho^2$ is
$\log(C\rho t^{-3/2})-3/2\geq\log(C\rho^{-2})-3/2>0$.
Its maximum on this interval is therefore at most
$C\rho^2(1+|\log\rho|)\leq C\rho$.
The interaction contribution is at most $4\|v\|_1\rho$.
This proves (6), uniformly even as $T\downarrow0$.
\end{proof}
